\begin{table}[htbp]
\centering
\caption{Bunching Estimates at Regional Help to Buy Price Caps}
\label{tab:bunching}
\begin{adjustbox}{max width=\textwidth}
\begin{tabular}{lrrrrrrr}
\hline\hline
Region & Cap (\pounds) & N & Bunching Ratio & SE & Placebo $b$ & Placebo SE & Excess Mass \\
\hline
\multicolumn{8}{l}{\textit{Panel A: New Builds (Post-Reform)}} \\
North East & 186,100 & 16,722 & 2.774*** & (0.393) & -0.080 & (0.135) & 208 \\
North West & 224,400 & 35,415 & 3.721*** & (0.334) & 2.331 & (0.118) & 436 \\
Yorkshire and The Humber & 228,100 & 27,922 & 3.827*** & (0.330) & -0.178 & (0.089) & 463 \\
West Midlands & 255,600 & 23,396 & 2.124*** & (0.320) & -1.517 & (0.082) & 192 \\
East Midlands & 261,900 & 32,283 & 1.843*** & (0.288) & -0.020 & (0.109) & 247 \\
South West & 349,000 & 29,634 & 0.579** & (0.283) & -0.575 & (0.105) & 56 \\
East of England & 407,400 & 33,018 & 0.280 & (0.262) & -1.765 & (0.103) & 28 \\
South East & 437,600 & 45,795 & 2.364*** & (0.294) & -1.166 & (0.086) & 299 \\
London & 600,000 & 29,450 & 3.055*** & (0.585) & -0.038 & (0.188) & 106 \\
\hline\hline
\end{tabular}
\end{adjustbox}
\begin{minipage}{\textwidth}
\footnotesize \textit{Notes:} Bunching ratio $b$ measures excess mass below the price cap relative to the polynomial counterfactual. Bootstrapped standard errors (500 iterations) in parentheses. Placebo $b$ is the bunching ratio for second-hand properties at the same cap (no bunching expected). $^{***}$, $^{**}$, $^{*}$ denote significance at the 1\%, 5\%, and 10\% levels.
\end{minipage}
\end{table}
